\chapter{Limitations}
\label{ch:limitations}

The items below set the scope of this project and the next steps that follow
from it.

\begin{itemize}
  \item \textbf{Live validation scope.} The live deployment had one
    participant, the researcher, in an n-of-1 design~\cite{vohra2015cent} with
    fictional adherence behaviour. It validates end-to-end operation on real
    infrastructure, and the work makes no claim about clinical effectiveness,
    usability, or adherence outcomes. An ethically approved, multi-site study
    with real patients is the next step, together with a formal clinician
    usability study of the dashboard.
  \item \textbf{Gold-set construction.} The 207 classifier inputs are
    synthetic, and one researcher authored and labeled them, so the classifier
    results describe this gold set. A second independent rater and an
    evaluation on collected patient replies would add inter-rater reliability
    and real-language coverage. Every adherence signal is also patient
    self-report, which varies with phrasing, recall window, and response
    format~\cite{stirratt2015selfreport}.
  \item \textbf{Safety-layer coverage.} Life-safety recall was 1.000 on the
    tested synthetic cases and across five repeated runs. That result bounds
    behaviour on tested cases, and deployed messages can differ from a
    benchmark~\cite{bedi2025testing}. Broader coverage of euphemisms,
    languages, spelling variants, and multi-symptom messages is future work.
    The layer routes to emergency services and does not diagnose or triage,
    and MedAdhere is not a crisis service.
  \item \textbf{Conversation evaluation.} Nine of the eighteen conversation
    metrics are scored by external semantic review against the hashed
    codebook, and the semantic core targets English. Scoring those metrics
    with independent raters and extending the core to the corpus's other
    languages are the next evaluation steps.
  \item \textbf{Platform and data.} The prototype runs in one AWS region on a
    Single-AZ database, with one administrator, one carrier, and one model
    provider, which suits a research deployment. Clinical use would add
    multi-AZ availability, role-based access, and a data-residency review,
    because inference runs through the Anthropic API in the United States. The
    research export is pseudonymized, not anonymous, so access control
    protects it until its deletion after grading. The adversarial gold-set
    slice samples prompt injection, a formalized attack
    space~\cite{liu2024promptinjection}, without covering it exhaustively.
\end{itemize}
